\documentclass{article}

% For figures
\usepackage{graphicx} % more modern
\usepackage{subfigure} 
\usepackage{tikz}
%\usepackage{pgfplots}
\usepackage{amsmath}

% For citations
%\usepackage[round]{natbib}
%\usepackage[backend=biber,style=alphabetic,citestyle=authoryear-comp]{biblatex}
\usepackage[authordate,bibencoding=auto,strict,backend=biber,natbib]{biblatex-chicago}
 
\addbibresource{abm.bib} %Imports bibliography file

% For algorithms
\usepackage{algorithm}
\usepackage{algorithmic}
\usepackage{hyperref}
\newcommand{\theHalgorithm}{\arabic{algorithm}}
\usepackage{mathtools}
\usepackage{amsmath}
\usepackage{kbordermatrix} %http://www.hss.caltech.edu/~kcb/TeX/kbordermatrix.sty

%BH additions for pagestyles
\usepackage{lastpage}
\usepackage{fancyhdr}
\pagestyle{fancy}
\fancyhf{}
\lhead{Haney, \textit{Beyond GDP}}
\rhead{Working Draft: Please do not distribute.}
\rfoot{Page \thepage \hspace{1pt} of \pageref{LastPage}}
\usepackage{multicol}
\renewcommand{\headrulewidth}{.25pt}


%%%%%%%%%%%%%%%%%%%%%%%%%%%%%%%%%%%%%%%%%%%%%%%%%%%%%%%%%%%%%%%%%%%%%%%%%%%%%%%%%%%%%%%%%%%%%%%%%%%%%%%%%%%%%%%%%%
\begin{document} 

\title{\textbf{Working Draft: Please do not distribute.}\\Beyond GDP: Measuring Economic Growth and Resilience
Using Ecological Network Analysis}
\author{Becky Roselius Haney~\footnote{Associate 
Professor of Economics, Calvin College, 1740 Knollcrest Cir SE, Grand Rapids, MI USA.
bhaney@calvin.edu.}}

\maketitle
\thispagestyle{empty} 

\begin{abstract} 
GDP is deceptively simplistic. Underneath this number
is a complex network of flows of goods, services, materials, and energy 
across
a large number of industries. Flows
of inputs to and outputs from each industry are used 
to create economic input-output (I-O) tables, from which
GDP is calculated.
These data have been meticulously collected
since the early 20th century, ever since Wasily Leontief developed
the input-output model of the economy. He later won the Nobel
prize in economics for this work.
These annual snapshots of the economic network 
provide insight into how the economy changes,
adapts, or evolves
over time. Many have proposed that economy may indeed evolve,
in a manner similar to an ecosystem. (See e.g. \cite{nelson_evolutionary_1985,arthur_complexity_1999,durlauf_complexity_2005, colander_complexity_2014})
With that in mind, ecologist \citet{ulanowicz_quantifying_2009} has
re-framed his metrics for ecological input-output tables
to use economic input-output tables. These metrics calculate 
measures of ``fitness for
evolution,'' or robustness of a system. 
Robust systems
have an optimal balance of what
Ulanowicz terms ``systemic
efficiency'' and ``resiliency capacity.'' 
Following Ulanowicz,
others have also applied 
these metrics  across a wide variety of economic settings. One
interesting finding so far is that  measures of economic robustness are
consistently
lower than those of ecosystem robustness. Does
this indicates that economies are
 more vulnerable to shocks than ecosystems? Or, is
the optimal balance of efficiency and capacity 
different for economies than for ecosystems? This paper 
sheds light on these questions by generating
Ulanowicz' metrics for a wide spectrum of simulated economies.
Using agent-based modelling, economies are generated that have varying
degrees of
industrialization and network structures. Their responses to
exogenous shocks are one measure of robustness. These
measures are compared to the ecological measures of robustness
to explore appropriate values of robustness for
economies as measured by ENA. The results suggest that... 
 

\end{abstract} 



\section{Introduction: Preventing the Dust Bowl}

Like the stock market crash of 1929, the massive ecological disaster 
that struck the Great Plains in the 1930's was largely unexpected.
The relatively cheap cost of land had enticed 
droves of settlers to move west, stake a claim,
and begin cultivating  acres of  prairie land. 
Agricultural productivity increased
over three hundredfold between the 1880
and 1930. Thriving farms and villages dotting the
previously barren prairie landscape foretold nothing
but a healthy, growing economic future.

However,  the rapid economic growth hid the fact that the
agricultural productivity and economic growth
of the Great Plains were coming at the expense of ecological 
resilience. Tilled acres and grazing lands gradually replaced
vast acres of what had seemed to be empty prairie land. 
However, the prairies had not been empty. They were dense ecosystems that had
been honed over time to thrive in the windy, arid desert climate of the
Great Plains. The sparse clumps 
of buffalo grass that gave the prairie landscape its 
deceptively barren look, had served an invaluable niche
role
in the prairie ecosystem. Over time, the grass had
evolved long roots to anchor the soil, and a 
cellular structure that could provide nutrition to the
mammals living in the arid conditions. Uprooting
the buffalo grass to make way for productive agricultural land 
seemed innocuous enough. And for several years, nothing seemed
to indicate differently. Inevitably, though, the drought years
returned and the devastating vulnerability
of the land was
exposed. 

From the early 1930s until the drought ended in 1939,
dust storms such as the one depicted
in Figure~\ref{fig:dust}  blew across the Great Plains. 
These ``black blizzards'' removed millions of tons
of topsoil and scattered the fine silt. 
Great Plains soil was blown all the way to the eastern seaboard.
Ships in the Atlantic ocean, as far as 300 miles off shore
reported appearances of the fine silt.~\citep{kassass_drought_1987}

\citet{kassass_drought_1987}
\citet[p. 36]{kassass_drought_1987}


\begin{figure}[H]
\centering\
\includegraphics[width=\linewidth]{Images/noaaphoto1935.jpg}
\caption{Photograph of a dust storm approaching Stratford, Texas, 1935. Credit: U.S. 
National Oceanic and Atmospheric Administration, George E. Marsh Album
(www.photolib.noaa.gov/htmls/theb1365.htm).
}
\label{fig:dust}
\end{figure}
Federal regulations on farming practices, WPA irrigation projects in the 1930's
and the end of the drought in 1939 eventually brought the
land back into balance. It took at least three decades,
however, to
return the land to a productive and livable state. 


The feast to famine gyrations of the desert southwest economy 
was not necessarily inevitable. In 1879, the first U.S. Geological Survey
was commissioned to map the more than 1.2 billion acres of 
the largely arid land owned by
the federal government, before settlers from the east 
began settling the Great Plains in earnest. 
In the ``Report on the Lands
of the Arid Regions of the United States,'' geologist
John Wesley Powell laid out an
ecologically-based plan for humans settlements.
Homestead boundaries would be based on 
the watershed and small systems of dams built, so that 
local communities would have property rights to their own water.

Powell's recommendations would not have created 
an environment conducive to large concentrations
of human population and farms in the desert southwest.
Instead, it would have been an loosely connected
network of smaller, water-based local communities.
Each community would self-regulate and maintain 
a scale appropriate to the water supply.
Instead, little planning and oversight accompanied
the great westward expansion. Homestead boundaries were 
drawn along neat, square grids. Access to water
involved massive water capture and redirection
projects, such as the Hoover and other large
dams built during this time. 
The Colorado river served as the most water
supply for the western states. Sharing this
one large water source involved negotiations and 
long-term commitments with distant
up-stream and down-stream communities.  
The 1922 Colorado river compact, which is
largely still in force today,
gives approximately 50\% of the approximately 600 cubic
meters per second water flow to California and
Colorado, and splits the rest between five other western states.
This water-sharing compact made it possible
for the massive Hoover Dam project
to collect and direct massive
amounts of water and energy 
to create cities in the desert.

The Hoover dam and cities in the desert such
Phoenix
would not have made sense under
Powell's ecological approach. 
Instead of large mega-cities such as
Los Angelas and Phoenix, located in 
the desert and heavily
dependent on an external water supply.
In many ways, the resulting network of large cities 
connected to massive public
infrastructure are efficient in that they 
take advantage of
agglomeration efficiencies and increasing
returns to scale. However, similar to the
dust bowl phenomenon, this type of
economic network can be vulnerable to water supply
shocks. 
This stands is contrast to the structure of the 
economy envisioned by Powell. The 
very large metropolises such as Phoenix
and L.A. would not exist. Instead, the economy
would be made up of many more smaller-sized communities,
scaled in proportion to the availability of 
water and power that could be supplied by
nearby series of small dams. The economic scale
and local management might have proven
more resilient  during 
the drought years of the dust bowl. The growth
and development of the economy over time might
have prevented
the contemporary water crises and water rights'
disputes between the western states, and between
the U.S. and Mexico. 
Indeed, a U.N report uses the development of the
Great Plains leading up to the dust bowl era as an
example of how not to undertake economic development.~\cite{kassass_drought_1987}.

The feast and famine gyrations of the Great Plains economy
suggests
that it could be fruitful to explore measures that 
can describe aspects of an economy beyond productivity
growth. Recently, \cite{arthur_complexity_1999} and 
others have argued that the economy grows, or evolves,
like an ecosystem. That is,
the economy is not well-described in aggregates of capital and labor,
but rather in the network of flows between heterogeneous
economic actors such as
industries, consumers, and the physical world. In other words,
the economy evolves like a complex adaptive system. 

The health of a complex adaptive system is measured 
differently than that of a machine. While the
total amount of output of a CAS is one measure, its
ability to innovate and re-create in the face of environmental
changes is equally important. 
If this is true, what metrics  
beyond GDP might be useful to evaluate the health of an
economy that evolves more akin to a garden than a machine?
The next section discusses the recent work of Ulanowicz, King, 
Khazzari and others who have applied ecological information metrics
used to understand the efficiency, resilience, and 
evolutionary robustness attributes of ecosystems
 to describe an economy or economic sector.


\section{Economies as Complex Adaptive Systems}

What are the root causes of economic growth and development? Indeed,
what exactly is economic growth and how does it differ from
economic development? What contributes to the resiliency of an economic system?



How might the efficiency, resiliency, and robust-ness
of a growing economy be described? What would an ``optimal balance of productivity
and resiliency'' look like? This is a question that biologists 
have considered in the
context of ecosystems. A few decades ago, ecologists
began using an input-output model to describe
ecosystems as primarily a
a network of flows, similar to an economy. Developed 
originally by Wasily Leontief (who won the Nobel prize for this work
in 1974), input-output tables developed the input-output model of the economy
are a modeling technique developed by 
ecosystem is modeled as
a network of flows, say of carbon that passes from prey
to predator when consumed, metrics developed
in information theory can measure efficiency and 
resilience and the optimal combination to create a
robust system. Ecological information metrics
provide additional insight regarding the health
of an ecosystem, beyond simply the quantity of 
flora and fauna that it supports. Ecosystems
that provide redundant forms of prey,
are more robust to a shock to the population of one of 
the sources of sustenance for predators. 

One promising source is the work of 
Ulanowicz~\cite{ulanowicz_quantifying_2009}
based on the work of MacKay~\cite{mackay_information_1969}
and others who have applied measures of ecosystem vitality to
economic vitality. Recently, the work of  Kharrazi and co-authors~\cite{kharrazi_examining_2015}\cite{kharrazi_quantifying_2013}\cite{,kharrazi_measuring_2016,kharrazi_network_2017},
King~\cite{king_information_2016} and others~\cite{huang_ecological_2014,zhang_analysis_2014,wang_inter-country_2016} have
applied these ecological information metrics to
the economy.

Ecological information
metrics provide summary measures of a network of flows 
that characterize system attributes such as:
how much the network produces (total system throughput, $TST$),
how efficient the network is in directing flows of 
inputs to their most productive uses, how
resilient the system is to shocks, and how
robust the system for the long-term by
an appropriate balance of efficiency and resiliency.

The simplified
ecosystem food chain in Figure~\ref{fig:ulano_ex}
can be used to illustrate how measures of
efficiency and resiliency and robustness
information metrics reflect characteristics of
a network.

\begin{figure}[H]
\centering\
\includegraphics[width=.9\linewidth]{Images/Ulano_ex.png}
\caption{Simplified network of carbon flows (in units of $mgCm^{-2}y^{-1}$) 
in a southern Florida wetland. 
\textit{Source: Author's calculations and Figure 1 from~\cite{ulanowicz_quantifying_2009}.}}
\label{fig:ulano_ex}
\end{figure}

One total system throughput (the sum of all flows)

To use ecological information metrics,  the network of flows are 
converted to an input-output
matrix format. The matrix contains a row and a column for each ``node,''
or species in the above ecosystem. In the $r$ x $c$ matrix $T$, 
cell $T_{ij}$
contains the volume of the flow from species \emph{i} to species \emph{j}.
For example, the above ecological network of flows could be
represented in the following input-output matrix:

\renewcommand{\kbldelim}{[}% Changes default Left  delimiter from ( to [
\renewcommand{\kbrdelim}{]}% Changes default Right delimiter from ( to [

\begin{equation*}
T = 
\kbordermatrix{
   & P & T & F & S & A          \cr
 P & 0 & 64.5 & 11.9 & 14.0 & 0 \cr
 T & 0 & 0 & 0 & 0 & 6.25 \cr
 F & 0 & 0 & 0 & 0 & 4.13 \cr
 S & 0 & 0 & 0 & 0 & 1.79 \cr
 A & 0 & 0 & 0 & 0 & 0 \cr}
 \end{equation*}
\\
\\
Biologist Ulanowicz and others have applied
information metrics to this ecosystem I-O
matrix. First, this ecosystem can be summarized
in terms of its efficient use of materials.
That is, how well do inputs flow to the species
with the ability to produce the most output with them. 
In the wetland ecosystem for example, notice
that for every unit of carbon flowing in, 
the Turtle species provides .10 units $(6.25/64.5)$
of carbon to
the upstream species (Alligators). Similarly, 
the Fish and
Snake species provide .35 and .13 units of 
carbon for the upstream species per 1 
unit of output.  Fish are the most efficient
producers of carbon. 

If Turtles and Snakes are much less efficient producers of carbon,
why have they not been winnowed out by the forces of 
natural selection? Because the long-term
sustainability of an ecosystem requires not only
efficient use of resources, but redundancy and
``slack'' capacity to allow recovery and
innovation as the environment changes. 

If the sole measure of the robustness of
an ecosystem were, 
all of the carbon would flow from Prawns through
Fish to the final user, Alligators,  as
in Figure~\ref{fig:ulano_ex1}. The
TST
of the system would
indeed grow from 102 units to 121.8 units.
However, the ability of the ecosystem to
withstand a shock to the Fish population
would drop to zero. 

\begin{figure}[H]
\centering\
\includegraphics[width=.9\linewidth]{Images/Ulano_ex1.png}
\caption{Simplified network of carbon flows (in units of $mgCm^{-2}y^{-1}$) 
in a southern Florida wetland. 
\textit{Source: Author's calculations and Figure 2 from~\cite{ulanowicz_quantifying_2009}.}}
\label{fig:ulano_ex1}
\end{figure}

The input-output model of inter-species flows
is based off the Leontief input-output matrix of
inter-industry flows for an economy.
Similar to an ecosystem, an economy consists of flows of
raw materials into industries, which transform
them into useful goods and services, which flow 
up the ``food chain'' to other industries
and ultimately to final users.
Economist Wassily Leontief won the Nobel memorial 
prize in Economics
for developing the input-output model 
of an economy. These ``I-O'' matrices are used extensively
by the U.S. Bureau of Economics Analysis (BEA)
and others to measure the output of industrial
sectors and the economy as a whole. I-O matrices
are also used to examine the inter-dependencies between 
industries. The I-O matrices can be used,
for example,
to forecast how a change in final demand will ripple 
backwards to affect employment levels in
all feeder industries. 

We come full circle then. Biology imported the I-O model 
from economics because of its usefulness to describe
ecosystem inter-species networks. We now explore
the usefulness of applying ecological
information metrics to summarize efficiency and
resiliency of the economy. The following section
uses the Florida wetland ecosystem example to illustrate
how it can be characterized  relative levels of efficiency and resiliency, an economy can be too.

%%%%%%%%%%%%%%%%%%%%%%%%%%%%%%%%%%%%%%%%%%%%%%%%%%%%%%%%%%%%%%%%%%%
\section{Describing Efficiency and Resiliency}
%%%%%%%%%%%%%%%%%%%%%%%%%%%%%%%%%%%%%%%%%%%%%%%%%%%%%%%%%%%%%%%%%%%

Ecological information metrics can provide summary measures for networks 
of flows. An ecosystem can be seen as a network of flows of 
carbon, lower members of the food chain provide energy in 
the form of carbon (?)
to those who consume them. These flows of carbon can be
represented as an input-output matrix, similar to economic
Leontief I-O matrices. Ecological information metrics are 
a subset of information metrics that measure the vitality
of an ecosystem, where vitality includes both efficiency and
resiliency. 

Like an ecosystem, an economy's robustness is also a combination
of efficiency and resiliency. Efficiency is a reflection of
using the most efficient pathways (industries) to convert
raw materials into useful goods and services. This is 
often equated with maximizing output, or Gross Domestic Product (GDP).

However, also like an ecosystem, an economy must maintain
some redundancy, or less efficient pathways to allow
flexibility to unforeseen shocks. For example, suppose the most
efficient primary input into electricity coal, in the sense that the 
kiloJoules of energy from \$1 worth of coal is greater than other
possible primary inputs. The most efficient economy would
funnel all of its resources toward making electricity from coal.
However, should a shock to the coal supply arise, the economy
would crash and require a significant recovery time to 
invent new electricity production methods and develop
the appropriate infrastructure.

Thus, to be robust, an economy needs to build in
at least some redundancy along with efficiency. Redundancy
includes not only alternative ways to produce key goods and
services, but also innovation. For technology to evolve,
at least some human and physical capital must be diverted from
the production line to work in the research lab. This reduces
efficiency, but increases the overall robustness of the economy.

How might the robustness of an economy be measured? 

%The following picture is copied and modified from the code in the manual at this website: https://www.bu.edu/math/files/2013/08/tikzpgfmanual.pdf (page 162)
\begin{center}
\begin{tikzpicture}[scale=6]
\draw[->] (-0.1,0) -- (1.1,0) node[right] {$\alpha=f(\textit{eff},\textit{res})$};
\draw[->] (0,-0.1) -- (0,0.5) node[left] {$F=-\alpha log(\alpha)$};
\draw[color=blue][domain=0.001:1] plot[id=F] function{-x*log(x)};
\draw (0.36 cm,1pt) -- (0.36,-1pt) node[below] {$\alpha^{*}$};
\draw (0.3 cm,.37) -- (0.44,.37) node[right] {$F'=0$};
% \foreach \y in {1,2,3}
%\draw (1pt,\y cm) -- (-1pt,\y cm) node[left] {$\y$};
\end{tikzpicture}
\end{center}

\begin{center}
\begin{tikzpicture}[scale=6]
\draw[->] (-0.1,0) -- (1.1,0) node[right] {$\alpha=f(\textit{eff},\textit{res})$};
\draw[->] (0,-0.1) -- (0,0.5) node[left] {$F=-\alpha log(\alpha)$};
\draw[color=blue][domain=0.001:1] plot[id=F] function{-x*log(x)};
\draw (0.36 cm,1pt) -- (0.36,-1pt) node[below] {$\alpha^{*}$};
\draw (0.3 cm,.37) -- (0.44,.37) node[right] {       $F'=0$};
\draw [color=red] (.2,.32188) circle (.2pt) node[above] {$A$};
\draw [color=red] (1,0) circle (.2pt) node[below] {$B$};
% \foreach \y in {1,2,3}
%\draw (1pt,\y cm) -- (-1pt,\y cm) node[left] {$\y$};
\end{tikzpicture}
\end{center}

\begin{center}
\begin{tikzpicture}[scale=6]
\draw[->] (-0.1,0) -- (1.1,0) node[right] {$\alpha=f(\textit{eff},\textit{res})$};
\draw[->] (0,-0.1) -- (0,0.5) node[left] {$F=-\alpha^{\beta} log(\alpha^\beta)$};
\draw[color=blue][domain=0.001:1] plot[id=F] function{-x**1.288*log(x**1.288)} ;
\draw[color=green][domain=0.001:1] plot[id=F] function{-x**2.7*log(x**2.7)} ;
\draw (.6,.2) node[right] [scale=.75] {$\beta=1.288$}
\draw (.8,.4) node[right] [scale=.75] {$\beta=\textit{e}$}

\end{tikzpicture}
\end{center}

\section{Application to the U.S. Energy Sector}
Similar to this previous work, we will apply ecological
information metric to energy sectors of the economy using
Leontief Input-Output data publicly available from the U.S.
Bureau of Economic Analysis and the U.S. Energy
Information Administration.

%\end{multicols}
\subsection{Ecological Information Metrics}

\begin{table}[H]
\caption{Ecological Information Metrics.}
\label{IOmetrics}
\vskip 0.15in
\begin{center}
\begin{sc}
\begin{tabular}{ l l}
\hline

Metric &  Formula  \\
\hline
&  \\
Systemic Efficiency
  & $SE = \displaystyle\sum_{ij}{T_{ij}} log(\frac{T_{ij}TST}{T_{i.}T_{.j}}) $ 
  \\[.25in]

Reserve Capacity 
  & $RC = - \displaystyle\sum_{ij}{T_{ij}} log(\frac{T_{ij}^2}{T_{i.}T_{.j}}) $ 
 \\[.25in]
  
Capacity (SE + RC)
  & $C = -\displaystyle\sum_{ij}{T_{ij}} log(\frac{T_{ij}}{TST}) $ 
   \\[.25in]

Organized Power  & $\alpha=\frac{SE}{SE + RC}$  \\[.25in]

Fitness
 
 & $F = -[\frac{e}{log(e)}] 
    \alpha^\beta log(\alpha^\beta)$ \\[.25in]

& $F' = -e \beta \alpha^{\beta - 1}
\bigg(
    \frac{log(\alpha^\beta)}  {log(e)} 
        + 1 
\bigg)$         \\[.25in]

Robustness & $R = TST \cdot F$\\[.25in]

& $\frac{\partial R}{\partial T_{ij}} = F + \frac{TST \cdot F'}{C}
\bigg(log \big(\frac{T_{ij} \cdot TST}{T_{i.} T_{.j}} \big) 
+ \alpha log \big( \frac{T^{2}_{ij}}{T_{i.}T_{.j}} 
\big) \bigg)$ \\[.25in]

First-order Conditions & $F = 1, F'=0$ \\[.25in]

\end{tabular}
\end{sc}
\end{center}
\vskip -0.1in
\end{table}

where \\
\begin{tabular}{r l}
$T_{ij}$  & quantity that flows from sector $i$ to sector $j$ \\[.1in]
$T_{i.}$ & sum of all flows from sector $i$ \\
$T_{.j}$ & sum of all flows into sector $j$ \\
$TST$ & $1$x$1$ scalar containing Total System Throughput (i.e. $T_{..})$\\
$\mathbf{IO}$ & $r$ x $c$ IO matrix of flows \\
$\mathbf{R}$ &  $r$ x $1$ vector of row sums of the IO matrix\\
$\mathbf{C}$ &  $1$ x $c$ vector of column sums of the IO matrix \\
$\beta$ &  parameter to allow optimal $\alpha$ 
    to vary across contexts  \\
$\circ$ & element-wise multiplication of two equally-dimensioned
matrices\\
$log(\mathbf{M})$ & element-wise log transformation of the 
entries in matrix $\mathbf{M}$\\
$\frac{1}{\mathbf{M}}$ & element-wise inverse of the entries in 
matrix $\mathbf{M}$ \\
$\mathbf{M}^2$ & element-wise square transformation of the 
entries in matrix $\mathbf{M}$ \\
$GRAND \sum \mathbf{M}$ & sum of all entries 
in matrix $\mathbf{M}$ (e.g. $TST=GRAND\sum \mathbf{IO})$ \\


\end{tabular}


\medskip

\subsection{Interpretation}

  The constant, $-\frac{e}{log(e)}$, in the formula for \emph{Fitness} 
  is constructed so that the function takes on the value of 1 at the maximum. 
  The adjustable parameter, $\beta$, 
  allows the optimal $alpha$ to vary across contexts. For example, 
  \cite{ulanowicz_quantifying_2009} suggests setting $\beta=1.288$ and 
  thus optimal $\alpha = .46$ for ecosystems. They estimated the 
  value based on observations of viable ecosystems and . By default, if $\beta=1$, then the Fitness function 
 is maximized at $\alpha = \frac{1}{e} = 0.3678$, an arbitrary choice.





\section*{Mathematical Appendix}
\begin{table}[H]
\caption{Matrix Formulation of Ecological Information Metrics.}
\label{IOmetrics_matrix}
\vskip 0.15in
\begin{sc}
\begin{center}
\begin{tabular}{ l l l }
\hline

Metric &  Formula & Matrix Formula \\
\hline
& & \\
%\hspace{.05in}
Systemic Efficiency
  & $SE = \displaystyle\sum_{ij}{T_{ij}}
  log(\frac{T_{ij}TST}{T_{i.}T_{.j}}) $ 
  & grand $\displaystyle\sum  \mathbf{IO} \circ 
  log\bigg(\mathbf{IO} \circ  \frac{TST}{\mathbf{R}\mathbf{C}}
  \bigg)$ \\[.25in]
  
  \multicolumn{3}{l}{\hspace{.5in} \texttt{mata code:
  SE=sum(IO:*(log(IO:*(TST:*(1:/(rowsum(IO)*colsum(IO)))))}} \\[.5in]

Reserve Capacity 
  & $RC = - \displaystyle\sum_{ij}{T_{ij}} log(\frac{T_{ij}^2}{T_{i.}T_{.j}}) $ 
  & -grand $\displaystyle\sum  \mathbf{IO} \circ 
  log\bigg(\mathbf{IO^2} \circ  \frac{1}{\mathbf{R} \mathbf{C}} \bigg)$ \\[.25in]
  
\multicolumn{3}{l}{\hspace{.5in} \texttt{mata code:
  RC=-sum(IO:*(log((IO:*IO):*(1:/(rowsum(IO)*colsum(IO)))))}} \\[.5in]
  
Capacity (SE + RC)
  & $C = -\displaystyle\sum_{ij}{T_{ij}} log(\frac{T_{ij}}{TST}) $ 
  & -grand $\displaystyle\sum  \mathbf{IO} \circ 
  log\bigg(\frac{1}{TST} \mathbf{IO}\bigg)$ \\[.25in]
  
\multicolumn{3}{l}{\hspace{.5in} \texttt{mata code:
  C=-sum(IO:*(log(IO:/TST)))}} \\[.5in]
  
\end{tabular}
\end{center}
\end{sc}
\vskip -0.1in
\end{table} 

%\bibliographystyle{plainnat}
%\bibliography{agent_based_modelling}

\printbibliography
\end{document} 


% For later...

\renewcommand{\abstractname}{Acknowledgements}
\begin{abstract}
This research project was generously underwritten by the Calvin College 
sabbatical program.
The computer simulations for this research were
performed using the Calvin College supercomputing cluster,
which was upgraded with funds provided under NSF Award \#1726260. 
The author thanks research assistants, Anthony Ditta and
Jiaming Jiang, for their diligent work during the original
programming and development of the agent-based
model used to produce the simulations for this research project. 
Discussions
with Matthew Kuperus Heun, Jeff Smith, Carey King,
Loren Haarsma, and John Sherwood contributed
greatly to this project. All errors are solely 
the responsibility of the author.
\end{abstract}
